% concatenated from g16New.tex
%g16New.tex

\documentclass{amsart}

\usepackage{amsmath,amssymb,amsfonts,mathtools,mathrsfs}
\usepackage{tikz-cd}
\usepackage{microtype}
\usepackage[colorlinks=true,linkcolor=blue,citecolor=blue,urlcolor=blue]{hyperref}

\newtheorem{theorem}{Theorem}[section]
\newtheorem{lemma}[theorem]{Lemma}
\newtheorem{proposition}[theorem]{Proposition}
\newtheorem{corollary}[theorem]{Corollary}
\theoremstyle{definition}
\newtheorem{definition}[theorem]{Definition}
\newtheorem{example}[theorem]{Example}
\newtheorem{question}[theorem]{Question}
\newtheorem{remark}[theorem]{Remark}

\newcommand{\id}{\operatorname{id}}
\newcommand{\End}{\operatorname{End}}
\newcommand{\Fun}{\operatorname{Fun}}
\newcommand{\Irr}{\operatorname{Irr}}
\newcommand{\Hom}{\operatorname{Hom}}
\newcommand{\Aut}{\operatorname{Aut}}
\newcommand{\Rep}{\operatorname{Rep}}
\newcommand{\Corep}{\operatorname{Corep}}
\newcommand{\Vect}{\operatorname{Vect}}
\newcommand{\Rad}{\operatorname{Rad}}
\newcommand{\rank}{\operatorname{rank}}
\newcommand{\Map}{\operatorname{Map}}
\newcommand{\Ind}{\operatorname{Ind}}
\newcommand{\Ocal}{\mathscr{O}}
\newcommand{\C}{\mathscr{C}}
\newcommand{\M}{\mathscr{M}}
\newcommand{\N}{\mathscr{N}}
\newcommand{\V}{\mathscr{V}}
\newcommand{\boxt}{\boxtimes}
\newcommand{\ot}{\otimes}
\newcommand{\Ctimes}{\mathbb{C}^{\times}}
\newcommand{\Gm}{\mathbb{G}_m}

\numberwithin{equation}{section}

\begin{document}

\title[Hopf $2$-cocycles for certain reductive groups]
{Hopf $2$-cocycles for certain affine algebraic reductive groups}

\author{Shlomo Gelaki}
\address{Department of Mathematics, Iowa State University, Ames, IA 50011, USA}
\email{gelaki@iastate.edu}

\date{August 27, 2026}

\keywords{affine algebraic reductive groups; semisimple tensor categories;
equivariantization; fiber functors; Hopf $2$-cocycles}

\begin{abstract}
Let $G$ be an affine algebraic reductive group over $\mathbb{C}$ whose identity
component is a torus $T$, and let $K:=G/T$.  We realize $\Rep(G)$ as an
equivariantization $\Rep(T)^K$ and describe its finite indecomposable semisimple
module categories in terms of equivariant module-category data over $\Rep(T)$.
For such an equivariantization, rank one is characterized by transitivity of the
induced action on the simple objects of the underlying $\Rep(T)$-module category,
and nondegeneracy of a projective cocycle on a point stabilizer.  We use this
criterion to parametrize fiber functors on $\Rep(G)$, prove that every such
fiber functor is classical, hence arises from a Hopf $2$-cocycle on
$\mathscr{O}(G)$, and classify (minimal) Hopf $2$-cocycle on
$\mathscr{O}(G)$. For commutative direct products $G=T\times K$, we also give the
canonical K\"unneth decomposition of the group of gauge classes, including the
mixed component, and compare it with our classification.
\end{abstract}

\maketitle
\tableofcontents

\section{Introduction}\label{sec:introduction}

Let $G$ be an affine algebraic group over $\mathbb{C}$, and let $\Rep(G)$ denote
the tensor category of finite-dimensional rational representations of $G$.
A Hopf $2$-cocycle on $\mathscr{O}(G)$ is equivalently a tensor structure on the
usual forgetful functor
\[
 F:\Rep(G)\longrightarrow \Vect.
\]
We call a tensor functor whose underlying linear functor is isomorphic to $F$ a
{\bf classical fiber functor}.

Hopf $2$-cocycles for finite groups were classified by Movshev \cite{Mo}, and the
nilpotent affine algebraic case was treated in \cite{EG1,EG2,G1}.  A general
classification for affine algebraic reductive groups remains open.  One of the
motivating questions of this paper is the following:

\begin{question}\cite[Question 7.2]{EG1}\label{q:torus}
If an affine algebraic reductive group admits a minimal Hopf $2$-cocycle, must its
identity component be a torus? \qed
\end{question}

We study the class of reductive groups for which the identity component is already
a torus.  Thus, throughout the main part of the paper, we have an exact sequence of affine algebraic groups 
\[
 1\longrightarrow T\longrightarrow G\longrightarrow K\longrightarrow 1,
\]
where $T$ is a {\bf complex torus} and $K$ is a {\bf finite} group.

The first ingredient is a tensor equivalence
\[
 \Rep(G)\simeq \Rep(T)^K.
\]
The action of $K$ on $\Rep(T)$ remembers not only the conjugation action of $K$ on
$T$, but also the extension class of $G$: the latter enters through the coherence
isomorphisms of the categorical action.  The second ingredient is the description
of finite semisimple module categories over an equivariantization.  In the present
setting, these are obtained from $L$-equivariant module categories over $\Rep(T)$,
where $L\subseteq K$ is a subgroup.

The decisive point for {\bf rank-one} module categories is an orbit--stabilizer
criterion.  If a finite group $L$ acts equivariantly on a finite semisimple category
$\M$, then the simple objects of $\M^L$ are indexed by orbits in $\Irr(\M)$ together
with irreducible projective representations of the corresponding stabilizers.
Consequently, $\M^L$ has rank one precisely when the action on $\Irr(\M)$ is
transitive and the stabilizer cocycle is nondegenerate.  In particular, rank one
does {\bf not} force the underlying category $\M$ itself to have rank one.
This phenomenon is exactly what produces mixed Hopf $2$-cocycles.

Using this criterion, we parametrize all rank-one module categories over
$\Rep(G)$ by the full equivariant data
\[
 (S,\Psi,L,U,\tau,\mu),
\]
subject to explicit coherence equations, transitivity of $U$, and nondegeneracy of
a stabilizer cocycle, and prove that every corresponding fiber functor has the
same dimension function as $F$ (see Theorem \ref{thm:module-and-fiber-classification}).  Since $\Rep(G)$ is semisimple, its underlying
linear functor is therefore isomorphic to $F$, and the tensor structure yields a
Hopf $2$-cocycle on $\mathscr{O}(G)$. Thus, we deduce the classification of (minimal) Hopf $2$-cocycle on $\mathscr{O}(G)$ (see Corollary \ref{cor:hopf-cocycle-classification} and Theorem \ref{prop:factorization}).

For a {\bf commutative} direct product $G=T\times K$, the character group is
$\widehat T\times\widehat K$, and the classification can be expressed canonically
as
\[
 H^2(\widehat T\times\widehat K,\Ctimes)
 \cong
 H^2(\widehat T,\Ctimes)
 \times \Hom(\widehat T\otimes_{\mathbb Z}\widehat K,\Ctimes)
 \times H^2(\widehat K,\Ctimes)
\]
(see Theorem \ref{thm:kunneth} and Corollary 
\ref{cor:kunneth-equivariantization-minimality}). The middle factor is essential since it can pair the radicals of the two diagonal
blocks, so minimality is governed by the full alternating bicharacter rather than
by nondegeneracy on either factor separately.

\medskip
\noindent
\textbf{Acknowledgements.}
This work was supported by Simons Foundation Award 963288.

\section{Preliminaries}\label{sec:preliminaries}

We work over the field $\mathbb{C}$ of complex numbers. We use the standard conventions for Hopf algebras and tensor categories, and refer to \cite{A,EGNO,M,S} for any unexplained notion.

\subsection{\texorpdfstring{Hopf $2$-cocycles}{Hopf 2-cocycles}}\label{sec:hopf-cocycles}

Let $H$ be a Hopf algebra over $\mathbb{C}$.  A \emph{Hopf $2$-cocycle} on $H$ is
a convolution-invertible functional $J\in(H\otimes H)^*$ satisfying
\begin{align*}
 J(a_1b_1,c)J(a_2,b_2)&=J(a,b_1c_1)J(b_2,c_2),\\
 J(a,1)&=\varepsilon(a)=J(1,a)
\end{align*}
for all $a,b,c\in H$.  Here and below we use Sweedler notation.

The cocycle $J$ defines a new Hopf algebra $H^J$.  Its underlying coalgebra is $H$, while its
multiplication is
\[
 a\cdot_J b=J^{-1}(a_1,b_1)a_2b_2J(a_3,b_3).
\]
Equivalently, $J$ defines a tensor structure on the forgetful functor
$$F:\Corep(H)\to\Vect.$$
For right $H$-comodules $X$ and $Y$, with coactions
$$x\mapsto x_{(0)}\otimes x_{(1)},\quad y\mapsto y_{(0)}\otimes y_{(1),}$$ 
the tensor constraint is
\[
 J_{X,Y}(x\otimes y)
 =J^{-1}(x_{(1)},y_{(1)})x_{(0)}\otimes y_{(0)}.
\]

For $a\in(H^*)^\times$, define
\[
 J^a(h,k)
 =a(h_1k_1)J(h_2,k_2)a^{-1}(h_3)a^{-1}(k_3).
\]
Two Hopf $2$-cocycles on $H$ are {\bf gauge equivalent} if they are related in this way.
Gauge-equivalence classes are in bijection with tensor structures on $F$ up to
monoidal natural isomorphism.

\subsection{Cotriangular Hopf algebras}\label{sec:cotriangular}

A cotriangular structure on a Hopf algebra $H$ is a convolution-invertible
functional $R\in(H\otimes H)^*$ satisfying
\begin{align*}
 R^{-1}&=R_{21},\\
 R(a,bc)&=R(a_1,b)R(a_2,c),\\
 R(ab,c)&=R(b,c_1)R(a,c_2),\\
 R(a_1,b_1)b_2a_2&=a_1b_1R(a_2,b_2)
\end{align*}
for all $a,b,c\in H$. It is {\bf minimal} if $R$ is nondegenerate.  By \cite[Proposition 2.1]{G1}, every
cotriangular Hopf algebra has a unique minimal cotriangular Hopf algebra quotient.

If $J$ is a Hopf $2$-cocycle on a cotriangular Hopf algebra $(H,R)$, then
\[
 R^J:=J_{21}^{-1}RJ
\]
is a cotriangular structure on $H^J$.  The cocycle $J$ is called {\bf minimal} if
$R^J$ is nondegenerate.

\subsection{Action groupoids and twisted action-groupoid algebras}
\label{sec:twisted-action-groupoid}
Let a group $\Gamma$ act on a set $X$.  The {\bf action groupoid}
$\Gamma\ltimes X$ is the category whose objects are the elements of $X$ and whose
arrows are
\[
 [g,x]:x\longrightarrow gx,
 \qquad g\in\Gamma,\ x\in X.
\]
Thus, two arrows $[g,x']$ and $[h,x]$ are composable precisely when $x'=hx$, and
then
\[
 [g,hx]\circ[h,x]=[gh,x].
\]
The identity arrow at $x$ is $[1,x]$, and the inverse of $[g,x]$ is
$[g^{-1},gx]$.

A normalized $2$-cocycle on this groupoid is a function
\[
 \omega:\Gamma\times\Gamma\times X\longrightarrow\Ctimes
\]
such that
\begin{equation}\label{eq:groupoid-cocycle}
 \omega(g,h;rx)\omega(gh,r;x)
 =\omega(h,r;x)\omega(g,hr;x)
\end{equation}
for all $g,h,r\in\Gamma$ and $x\in X$, and
$\omega(1,g;x)=\omega(g,1;x)=1$.  The corresponding {\bf twisted action-groupoid
algebra} $\mathbb C^{\omega}[\Gamma\ltimes X]$ is the vector space of finitely
supported linear combinations of the arrows, with multiplication
\begin{equation}\label{eq:twisted-groupoid-product}
 [g,x'][h,x]
 =
 \begin{cases}
  \omega(g,h;x)[gh,x],&x'=hx,\\
  0,&x'\neq hx.
 \end{cases}
\end{equation}
Equation~\eqref{eq:groupoid-cocycle} is exactly the associativity condition for
\eqref{eq:twisted-groupoid-product}.  The elements
\[
 e_x:=[1,x],\qquad x\in X,
\]
are pairwise orthogonal idempotents; if $X$ is finite, their sum is the identity
of $\mathbb C^{\omega}[\Gamma\ltimes X]$.

\begin{lemma}[Reduction of a transitive action groupoid]
\label{lem:transitive-groupoid-reduction}
Assume that the action of $\Gamma$ on the finite set $X$ is transitive.  Fix
$x_0\in X$, let
\[
 H=\Gamma_{x_0}:=\{g\in\Gamma\mid gx_0=x_0\},
 \qquad
 \omega_0(g,h):=\omega(g,h;x_0),\quad g,h\in H,
\]
and set $\mathcal A:=\mathbb C^{\omega}[\Gamma\ltimes X]$.  Then
\[
 e_{x_0}\mathcal A e_{x_0}\cong\mathbb C^{\omega_0}[H].
\]
Moreover, $e_{x_0}$ is a full idempotent and, after choosing one arrow from
$x_0$ to every object of $X$, there is an algebra isomorphism
\begin{equation}\label{eq:transitive-groupoid-matrix-reduction}
 \mathcal A
 \cong
 \operatorname{Mat}_{|X|}\!\left(\mathbb C^{\omega_0}[H]\right).
\end{equation}
Consequently, restriction to the corner gives an equivalence
\[
 {\rm Mod}(\mathcal A)
 \simeq
 {\rm Mod}(\mathbb C^{\omega_0}[H]).
\]
\end{lemma}

\begin{proof}
The corner $e_{x_0}\mathcal A e_{x_0}$ is spanned by the loops $[h,x_0]$ with
$h\in H$, and their multiplication is
\[
 [g,x_0][h,x_0]=\omega_0(g,h)[gh,x_0].
\]
This proves the identification with $\mathbb C^{\omega_0}[H]$.

For each $x\in X$, choose $r_x\in\Gamma$ with $r_xx_0=x$, taking
$r_{x_0}=1$, and let
\[
 p_x:=[r_x,x_0].
\]
The two products
\[
 [r_x^{-1},x]p_x
 \quad\text{and}\quad
 p_x[r_x^{-1},x]
\]
are nonzero scalar multiples of $e_{x_0}$ and $e_x$, respectively.  The cocycle
identity shows that the two scalars are equal.  Hence one can rescale
$[r_x^{-1},x]$ to an arrow $q_x$ such that
\begin{equation}\label{eq:connecting-arrows}
 q_xp_x=e_{x_0},
 \qquad
 p_xq_x=e_x.
\end{equation}
In particular, $e_x=p_xe_{x_0}q_x$ belongs to
$\mathcal A e_{x_0}\mathcal A$ for every $x$, so $e_{x_0}$ is full.

Let $E_{x',x}$ denote the standard matrix unit indexed by $x',x\in X$.  For
$a\in e_{x'}\mathcal A e_x$, define
\begin{equation}\label{eq:explicit-groupoid-matrix-map}
 \Phi(a)=E_{x',x}\otimes q_{x'}ap_x.
\end{equation}
The middle factor belongs to $e_{x_0}\mathcal A e_{x_0}$.  Using
\eqref{eq:connecting-arrows}, one checks that
\[
 \Phi(ab)=\Phi(a)\Phi(b)
\]
whenever $a$ and $b$ are composable, while both sides vanish otherwise.  The
inverse map is
\[
 E_{x',x}\otimes b\longmapsto p_{x'}bq_x,
 \qquad
 b\in e_{x_0}\mathcal A e_{x_0}.
\]
This proves \eqref{eq:transitive-groupoid-matrix-reduction}; the stated Morita
equivalence follows, for example, from the mutually inverse functors
\[
 N\longmapsto e_{x_0}N,
 \qquad
 W\longmapsto \mathcal A e_{x_0}
       \otimes_{e_{x_0}\mathcal A e_{x_0}}W.
\]
\end{proof}

\subsection{Actions and equivariantizations}\label{sec:actions}

Let $K$ be a {\bf finite} group and $\C$ a {\bf semisimple} tensor category (not necessarily finite). Recall (see, e.g., \cite[Definition 2.7.1]{EGNO}) that an action of $K$
on $\C$ consists of tensor autoequivalences $A_k$ and tensor natural
isomorphisms
\[
 \gamma_{k,\ell}:A_kA_\ell\xrightarrow{\cong}A_{k\ell},
 \qquad k,\ell\in K,
\]
satisfying the unit and pentagon axioms.

Recall (see, e.g., \cite[Definition 2.7.2]{EGNO}) that a $K$-equivariant object is a pair $(X,u)$, where $X\in\C$ and
\[
 u_k:A_k(X)\xrightarrow{\cong}X
\]
satisfy
\begin{equation}\label{eq:equivariant-object}
 u_kA_k(u_\ell)=u_{k\ell}\gamma_{k,\ell}(X).
\end{equation}
The category of equivariant objects is the equivariantization $\C^K$.
For the trivial action on $\Vect$, one recovers $\Vect^K\simeq\Rep(K)$.

\subsection{Equivariant module categories}\label{sec:equivariant-modules}

Let $\M$ be a left $\C$-module category.  For $k\in K$, define the twisted module
category $\M^k$ to have the same underlying category as $\M$, and action
\begin{equation}\label{eq:module-twist}
 X\otimes^k M:=A_{k^{-1}}(X)\otimes M.
\end{equation}
With this convention, $k\mapsto\M^k$ is a left action on equivalence classes of
$\C$-module categories.  More precisely, the coherence $\gamma$ induces canonical
module equivalences
\[
 \Gamma_{k,\ell}:(\M^\ell)^k\xrightarrow{\simeq}\M^{k\ell};
\]
the module structure of $\Gamma_{k,\ell}$ is induced by
$\gamma_{\ell^{-1},k^{-1}}$.

Let $L\subseteq K$ be a subgroup.  Recall \cite{ENO} that an $L$-equivariant structure on $\M$ consists of module
equivalences
\[
 U_k:\M^k\xrightarrow{\simeq}\M,
 \qquad k\in L,
\]
and module natural isomorphisms
\begin{equation}\label{eq:mu-abstract}
 \mu_{k,\ell}:U_k\circ(U_\ell)^k
 \xrightarrow{\cong}
 U_{k\ell}\circ\Gamma_{k,\ell},
\end{equation}
subject to the usual unit and pentagon conditions.  We always use normalized data:
$U_1=\id$, and the unit components of $\mu$ are identities.

An equivariant object is a pair $(M,v)$ with
\[
 v_k:U_k(M)\xrightarrow{\cong}M
\]
satisfying
\begin{equation}\label{eq:equivariant-module-object}
 v_kU_k(v_\ell)=v_{k\ell}\mu_{k,\ell}(M).
\end{equation}
The resulting finite semisimple category is denoted $\M^L$.  Recall \cite{ENO} that it carries a natural
$\C^K$-module structure; on underlying objects the action is the original action
$X\otimes M$ in $\M$.

For the trivial action on $\Vect$, an equivariant structure is determined by a
normalized cocycle $\mu\in Z^2(L,\Ctimes)$, and
\[
 \Vect^L\simeq\Rep(L,\mu),
\]
where $\rho(k)\rho(\ell)=\mu(k,\ell)\rho(k\ell)$.

\subsection{Simple objects and rank}\label{sec:orbit-stabilizer}

Let $L$ act equivariantly on a finite semisimple category $\M$.  Then the functors $U_k$
permute $\Irr(\M)$.  For $M\in\Irr(\M)$, let
\[
 L_M:=\{k\in L\mid U_k(M)\cong M\};
\]
it is a subgroup of $L$, hence of $K$. After choosing isomorphisms $U_k(M)\cong M$ for $k\in L_M$, the coherence data
induces a well-defined cohomology class
$$[\alpha_M]\in H^2(L_M,\Ctimes).$$ 
We call a cocycle on a finite group
{\bf nondegenerate} when its twisted group algebra is simple.

We have the following orbit--stabilizer description.

\begin{theorem}\label{thm:orbit-stabilizer}
For each $L$-orbit $\mathcal O\subseteq\Irr(\M)$, choose $M\in\mathcal O$.  Then
\[
 \M^L\simeq
 \bigoplus_{\mathcal O\in L\backslash\Irr(\M)}
 \Rep(L_M,\alpha_M),
\]
as an abelian category. Consequently,
\[
 \rank(\M^L)=
 \sum_{\mathcal O\in L\backslash\Irr(\M)}
 \#\Irr\bigl(\mathbb{C}^{\alpha_M}[L_M]\bigr).
\]
In particular, $\M^L$ has rank one if and only if $L$ acts transitively on
$\Irr(\M)$ and the twisted group algebra $\mathbb{C}^{\alpha_M}[L_M]$ is simple for one,
equivalently every, $M\in\Irr(\M)$.
\end{theorem}

\begin{proof}
Fix an orbit $\mathcal O$, and choose a representative simple object $M_x$ for
each $x\in\mathcal O$.  Choose isomorphisms
\[
 c_{k,x}:U_k(M_x)\xrightarrow{\cong}M_{kx}.
\]
Since every $M_x$ is simple, comparison of the two isomorphisms from
$U_kU_\ell(M_x)$ to $M_{k\ell x}$ gives a unique scalar
$\omega(k,\ell;x)\in\Ctimes$, which we normalize by
\begin{equation}\label{eq:groupoid-cocycle-from-equivariance}
 c_{k,\ell x}\,U_k(c_{\ell,x})
 =\omega(k,\ell;x)^{-1}
   c_{k\ell,x}\,\mu_{k,\ell}(M_x).
\end{equation}
The pentagon for $\mu$ says precisely that $\omega$ satisfies
\eqref{eq:groupoid-cocycle}; thus $\omega$ is a normalized $2$-cocycle on the
action groupoid $L\ltimes\mathcal O$ of
\S\ref{sec:twisted-action-groupoid}.

Every object of $\M$ supported on $\mathcal O$ has the form
\[
 N=\bigoplus_{x\in\mathcal O}V_x\otimes M_x.
\]
Under the chosen maps $c_{k,x}$, an equivariant structure on $N$ is equivalent
to a family of linear maps
\[
 \rho(k,x):V_x\longrightarrow V_{kx}
\]
satisfying
\begin{equation}\label{eq:groupoid-representation-relation}
 \rho(k,\ell x)\rho(\ell,x)
 =\omega(k,\ell;x)\rho(k\ell,x).
\end{equation}
Relation~\eqref{eq:groupoid-representation-relation} says exactly that the arrow
$[k,x]$ acts from $V_x$ to $V_{kx}$, and that these arrow actions define a module
over $\mathbb C^{\omega}[L\ltimes\mathcal O]$.  Hence the full subcategory of
$\M^L$ supported on $\mathcal O$ is equivalent to
\[
 {\rm Mod}(\mathbb C^{\omega}[L\ltimes\mathcal O]).
\]

Choose $M=M_{x_0}$ in the orbit $\mathcal O$.  The isotropy group at $x_0$ is $L_M$, and the
restriction of $\omega$ to its loops represents the class $[\alpha_M]$.  By
Lemma~\ref{lem:transitive-groupoid-reduction},
\[
 {\rm Mod}(\mathbb C^{\omega}[L\ltimes\mathcal O])
 \simeq
 {\rm Mod}(\mathbb C^{\alpha_M}[L_M])
 =\Rep(L_M,\alpha_M).
\]
Taking the direct sum over the orbits proves the categorical decomposition.  The
rank formula follows by counting the simple modules of the twisted stabilizer
algebras, and the rank is one exactly under the two conditions stated in the
theorem.
\end{proof}

\begin{remark}\label{rem:no-rank-inequality}
There is no general inequality comparing $\rank(\M^L)$ and $\rank(\M)$.  For
example, if $\mathbb{Z}/2\mathbb{Z}$ interchanges the two simple summands of $\Vect\oplus\Vect$, then
$(\Vect\oplus\Vect)^{\mathbb{Z}/2\mathbb{Z}}\simeq\Vect$. \qed
\end{remark}

\subsection{Module categories over an equivariantization}\label{sec:module-classification-general}

Let $\C\rtimes K$ be the crossed-product tensor category (see, e.g., \cite[Definition 4.15.5]{EGNO}).  Recall that the category $\C$ is a
left $\C\rtimes K$-module category, and the standard duality construction gives a tensor equivalence
\[
 (\C\rtimes K)^*_{\C}=\Fun_{\C\rtimes K}(\C,\C)\simeq\C^K.
\]
Explicitly, an equivariant object $(X,u)$ acts on $\C$ by the module endofunctor
$X\otimes-$; the equivariant structure supplies the module constraints for the
homogeneous components of $\C\rtimes K$.  Evaluation at the tensor unit gives the
inverse construction. (For the fusion case see, e.g., \cite[Example 7.12.25]{EGNO}.)

\begin{proposition}\label{prop:equivariantization-modules}
Let $T$ be a complex torus, let $\C:=\Rep(T)$, and let a finite group $K$ act
on $\C$. Then every finite indecomposable semisimple $\C^K$-module category is
equivalent to $\M^L$ for some subgroup $L\subseteq K$ and some indecomposable
$L$-equivariant $\C$-module category $\M$.

Moreover, two such pairs determine equivalent $\C^K$-module categories precisely when they
are related by conjugation in $K$ and an equivariant $\C$-module equivalence.
\end{proposition}

\begin{proof}
By the preceding remarks, the usual evaluation and coevaluation functors 
identify finite semisimple module categories over the Morita-equivalent tensor
categories $\C^K$ and $\C\rtimes K$. (Note that no infinite direct sums occur since $K$ is finite and every module category
under consideration has finitely many simple objects.)

Thus, it remains to describe finite indecomposable module categories over
$\C\rtimes K$.  After restriction to the neutral component $\C$, such a module
category decomposes into finitely many indecomposable $\C$-module summands.  The
invertible homogeneous objects $\mathbf 1\boxtimes k$, $k\in K$, permute these summands.
Indecomposability over $\C\rtimes K$ implies that the permutation action is
transitive.  If $\M$ is one summand and $L$ is its stabilizer, the homogeneous
action and its associativity constraints give an $L$-equivariant structure on
$\M$.  Conversely, induction from this stabilizer reconstructs the
$\C\rtimes K$-module category, and under the preceding Morita equivalence it
corresponds to $\M^L$. (This is the argument of \cite[Proposition 5.4]{ENO}; the
finiteness observations above justify it in the present locally finite setting.)
\end{proof}

\section{The reductive case}\label{sec:reductive}

From now on, let $G$ be an affine algebraic reductive group whose identity
component is a torus $T$, and let $K:=G/T$.

\subsection{\texorpdfstring{The tensor equivalence $\Rep(G)\simeq\Rep(T)^K$}{The tensor equivalence for Rep(G)}}
\label{sec:tensor-equivalence}

Consider the exact sequence
\begin{equation}\label{eq:extension}
 1\rightarrow T\rightarrow G\xrightarrow{\pi}K\rightarrow 1.
\end{equation}
Choose a set-theoretic section $\sigma:K\xrightarrow{1:1} G$ with $\sigma(1)=1$.  Define
\[
 \theta_k(t):=\sigma(k)t\sigma(k)^{-1},
 \qquad
 \eta(k,\ell):=\sigma(k)\sigma(\ell)\sigma(k\ell)^{-1}
\]
for all $k,\ell\in K$ and $t\in T$.

\begin{lemma}\label{lem:extension-data}
The maps $\theta$ and $\eta$ have the following properties:
\begin{enumerate}
\item $\theta:K\to\Aut(T)$ is an action, independent of the chosen section.
\item $\eta\in Z^2(K,T)$ for this action; explicitly,
\[
 \theta_k(\eta(\ell,m))\eta(k,\ell m)
 =\eta(k,\ell)\eta(k\ell,m)
\]
for all $k,\ell,m\in K$. Its cohomology class is independent of $\sigma$.
\item 
If $N:=\{t\in T\mid t^{|K|}=1\}$, then $[\eta]$ lies in the image of
\[
 H^2(K,N)\longrightarrow H^2(K,T).
\]
After changing $\sigma$ by a $T$-valued $1$-cochain, one may therefore choose an
$N$-valued representative of the extension class.
\item 
The map $T\times K\to G$, $(t,k)\mapsto t\sigma(k)$, is an isomorphism of
varieties, and the multiplication is
\[
 (t,k)(s,\ell)=\bigl(t\theta_k(s)\eta(k,\ell),k\ell\bigr).
\]
\end{enumerate}
\end{lemma}

\begin{proof}
(1), (2), and (4) follow directly from associativity in $G$
and the commutativity of $T$.  

(3) Use the exact sequence of
$K$-modules
\[
 1\longrightarrow N\longrightarrow T
 \xrightarrow{t\mapsto t^{|K|}}T\longrightarrow1.
\]
The power map is surjective because $T$ is a complex torus, while multiplication by
$|K|$ annihilates $H^2(K,T)$.  Thus the long exact cohomology sequence gives the claim.
\end{proof}

Let $\widehat{T}$ be the {\bf group of characters} of $T$. Recall that $\Rep(T)$ is a semisimple pointed tensor category with simple objects 
$$\{\delta_x\mid x\in \widehat{T}\}.$$
For $x\in \widehat{T}$ and $k\in K$, define
\[
 x^k:=x\circ\theta_{k^{-1}}.
\]
For $X\in\Rep(T)$, let $A_k(X)$ be the pullback of $X$ along $\theta_{k^{-1}}$.
Thus $A_k(\delta_x)=\delta_{x^k}$.  Define a tensor natural isomorphism
\begin{equation}\label{eq:gamma-extension}
 \gamma_{k,\ell}(X)
 :=\rho_X\!\left(\theta_{(k\ell)^{-1}}(\eta(k,\ell))\right)
 :A_kA_\ell(X)\xrightarrow{\cong}A_{k\ell}(X).
\end{equation}
The cocycle identity for $\eta$ is exactly the pentagon identity for $\gamma$, so
$(A,\gamma)$ is an action of $K$ on $\Rep(T)$.

\begin{theorem}\label{thm:tensor-equivalence}
The functor
\[
 \Phi:\Rep(T)^K\longrightarrow\Rep(G),
 \qquad
 (X,u)\longmapsto\bigl[t\sigma(k)\mapsto\rho_X(t)u_k\bigr],
\]
is a tensor equivalence.
\end{theorem}

\begin{proof}
Because $u_k:A_k(X)\xrightarrow{\cong} X$ is $T$-linear,
\[
 u_k\rho_X(s)=\rho_X(\theta_k(s))u_k.
\]
Moreover, \eqref{eq:equivariant-object} and \eqref{eq:gamma-extension} give
\[
 u_ku_\ell=\rho_X(\eta(k,\ell))u_{k\ell}.
\]
Together with Lemma~\ref{lem:extension-data}(4), these identities show that the
displayed formula defines a representation of $G$.

Conversely, if $(X,\pi)$ is a $G$-representation, restrict it to $T$ and set
$u_k:=\pi(\sigma(k))$.  Then $u_k:A_k(X)\to X$ is $T$-linear, and
\[
 u_ku_\ell=\pi(\eta(k,\ell))u_{k\ell}
 =u_{k\ell}\gamma_{k,\ell}(X).
\]
Thus $(X|_T,u)$ is a $K$-equivariant object.  

Finally, it is straightforward to verify that the two constructions above are mutually
inverse and preserve tensor products.
\end{proof}

For completeness, we record the corresponding description of irreducible
representations of $G$.  For $x\in \widehat{T}$, let
\[
 K_x:=\{k\in K\mid x^k=x\},
 \qquad
 \alpha_x(k,\ell)=x(\eta(k,\ell)),\quad k,\ell\in K_x.
\]
The restriction of $\eta$ to $K_x$ evaluates to the scalar cocycle $\alpha_x$.

\begin{proposition}\label{prop:simple-G-modules}
Isomorphism classes of simple modules over $G$ are parametrized by $K$-orbits in $\widehat{T}$
together with simple $\alpha_x$-projective representations of $K_x$.  If $W$ is
such a projective representation, the corresponding simple $G$-module is
\[
 \Ind_{G_x}^{G}(W\otimes\delta_x),
 \qquad G_x=\pi^{-1}(K_x),
\]
where $t\sigma(k)\in G_x$ acts on $W\otimes\delta_x$ by
$x(t)\rho_W(k)$.
\end{proposition}

\begin{proof}
The formula defines a representation of $G_x$ because $x$ is $K_x$-invariant and
$\rho_W(k)\rho_W(\ell)=\alpha_x(k,\ell)\rho_W(k\ell)$.  The assertion now follows
from Theorem~\ref{thm:tensor-equivalence} and the orbit--stabilizer description of
simple objects in an equivariantization.
\end{proof}

\subsection{\texorpdfstring{Finite module categories over $\Rep(T)$}{Finite module categories over Rep(T)}}
\label{sec:module-cats-torus}

Let $S\subseteq \widehat{T}$ be a subgroup of {\bf finite} index.  A normalized
{\bf module $2$-cocycle} on the transitive $\widehat{T}$-set $\widehat{T}/S$ is a function
\[
 \Psi:\widehat{T}\times \widehat{T}\times (\widehat{T}/S)\longrightarrow\Ctimes,
 \qquad (x,z,y)\longmapsto\Psi(x,z;y),
\]
satisfying
\begin{equation}\label{eq:module-cocycle}
 \Psi(x,z;wy)\Psi(xz,w;y)
 =\Psi(z,w;y)\Psi(x,zw;y)
\end{equation}
for all $x,z,w\in \widehat{T}$ and $y\in \widehat{T}/S$, together with
$\Psi(1,x;y)=\Psi(x,1;y)=1$.

Define $\M(S,\Psi)$ to be the finite semisimple category with simple objects
$$\{\delta_y\mid y\in \widehat{T}/S\},$$
action
\[
 \delta_x\otimes\delta_y=\delta_{xy},
\]
and module associativity constraint
\[
 m_{x,z,y}=\Psi(x,z;y)\id_{\delta_{xzy}}.
\]
Equation \eqref{eq:module-cocycle} is precisely the module pentagon.

\begin{proposition}[\cite{G2}; see also \cite{EGNO}]\label{prop:module-cats-torus}
Every finite indecomposable semisimple $\Rep(T)$-module category is equivalent to
$\M(S,\Psi)$ for some finite-index subgroup $S\subseteq \widehat{T}$ and some normalized module
$2$-cocycle $\Psi$.  Its equivalence class depends only on the cohomology class of
$\Psi$.  By Shapiro's lemma, these cohomology classes are naturally identified with
$H^2(S,\Ctimes)$.
\end{proposition}

\begin{proof}
Each invertible simple object $\delta_x$ acts by an autoequivalence, and hence $\widehat{T}$
acts on the finite set of simple objects of the module category.  Indecomposability
is equivalent to transitivity of this action.  Choose a simple object $M_0$, let
$S\subseteq \widehat{T}$ be its stabilizer, and identify the simple objects with $\widehat{T}/S$.
After choosing normalized isomorphisms
\[
 c_{x,y}:\delta_x\otimes M_y\xrightarrow{\cong}M_{xy},
\]
the module associativity constraint becomes multiplication by scalars
$\Psi(x,z;y)$.  The module pentagon and unit axioms give
\eqref{eq:module-cocycle} and the stated normalizations.  

Conversely, those
identities define a module category.  Replacing the maps $c_{x,y}$ by scalar
multiples changes $\Psi$ by a module coboundary, and every such coboundary arises
this way.  

Finally, the coefficient module is the coinduced $\widehat{T}$-module
$$\Map(\widehat{T}/S,\Ctimes),$$
so Shapiro's lemma identifies its second cohomology with
$H^2(S,\Ctimes)$.
\end{proof}

\subsubsection{Explicit equivariant structures}\label{sec:explicit-equivariant-data}

Fix $\M(S,\Psi)$ and a subgroup $L\subseteq K$.  We use the action on $\widehat{T}$ introduced in
\S\ref{sec:tensor-equivalence}.  For $k\in L$, let
\[
 \Psi^k(x,z;y)=\Psi(x^{k^{-1}},z^{k^{-1}};y),
 \qquad
 \gamma_{\ell^{-1},k^{-1}}(\delta_x)
 =\bar\gamma_{k,\ell}(x)\,\id_{\delta_{x^{(k\ell)^{-1}}}}.
\]
Thus $\M(S,\Psi)^k$ has action
$\delta_x\otimes^k\delta_y=\delta_{x^{k^{-1}}y}$ and associator $\Psi^k$.

\begin{theorem}\label{thm:explicit-equivariant-data}
The category $\M(S,\Psi)$ admits an $L$-equivariant structure only if $S$ is
$L$-stable.  If $S$ is $L$-stable, normalized $L$-equivariant structures are
parametrized by data
\[
 U_k\in\operatorname{Sym}(\widehat{T}/S),
 \qquad
 \tau_k:\widehat{T}\times (\widehat{T}/S)\to\Ctimes,
 \qquad
 \mu_{k,\ell}:\widehat{T}/S\to\Ctimes,
\]
for $k,\ell\in L$, satisfying the following equations for all
$i,k,\ell\in L$, $x,z\in \widehat{T}$, and $y\in \widehat{T}/S$:
\begin{equation}
 U_1=\id,\quad
 U_kU_\ell=U_{k\ell},\quad
 U_k(xy)=x^kU_k(y),
 \label{eq:U-data}
 \end{equation}
 \begin{equation}
 \tau_1(x,y)=1,\quad
 \tau_k(1,y)=1,
 \label{eq:tau-normalized}
\end{equation}
\begin{equation}\label{eq:tau-coherence}
 \tau_k(z,y)\tau_k(x,z^{k^{-1}}y)\Psi^k(x,z;y)
 =\Psi(x,z;U_k(y))\tau_k(xz,y),
\end{equation}
\begin{equation}\label{eq:mu-module}
 \bar\gamma_{k,\ell}(x)\tau_{k\ell}(x,y)
 \mu_{k,\ell}(x^{(k\ell)^{-1}}y)
 =\mu_{k,\ell}(y)\tau_k(x,U_\ell(y))
 \tau_\ell(x^{k^{-1}},y),
 \end{equation}
\begin{align}
 \mu_{1,k}(y)&=\mu_{k,1}(y)=1,
 \label{eq:mu-normalized}\\
 \mu_{i,k}(U_\ell(y))\mu_{ik,\ell}(y)
 &=\mu_{k,\ell}(y)\mu_{i,k\ell}(y).
 \label{eq:mu-pentagon}
\end{align}
Here $U_k$ is the permutation induced by the underlying functor of
$$U_k:\M(S,\Psi)^k\to\M(S,\Psi),$$
while $\tau_k$ and $\mu_{k,\ell}$ are the scalar
components of its module structure and coherence isomorphisms.
\end{theorem}

\begin{proof}
A module equivalence $U_k:\M(S,\Psi)^k\xrightarrow{\simeq}\M(S,\Psi)$ permutes the simple objects.
Its module constraint has components
\[
 U_k(\delta_{x^{k^{-1}}y})\longrightarrow
 \delta_x\otimes U_k(\delta_y).
\]
Equality of the two simple labels is equivalent to
$U_k(xy)=x^kU_k(y)$, which also implies that $S$ is $L$-stable.  The module-functor
pentagon and unit axiom are exactly \eqref{eq:tau-normalized} and
\eqref{eq:tau-coherence}.

The natural transformation
\[
 \mu_{k,\ell}:U_k(U_\ell)^k
 \longrightarrow U_{k\ell}\Gamma_{k,\ell}
\]
can have a nonzero component only when $U_kU_\ell=U_{k\ell}$.  The condition that
$\mu_{k,\ell}$ be a module natural transformation gives
\eqref{eq:mu-module}; the appearance of $\bar\gamma_{k,\ell}$ comes from the module
structure of $\Gamma_{k,\ell}$.  Finally, its pentagon is
\eqref{eq:mu-pentagon}.  

Conversely, the displayed equations reconstruct all the
required module functors and coherence maps.
\end{proof}

\subsubsection{Simple objects and the rank-one criterion}
\label{sec:rank-one-criterion}

For $y\in \widehat{T}/S$, let
\[
 L_y:=\{k\in L\mid U_k(y)=y\},
 \qquad
 \mu_y(k,\ell)=\mu_{k,\ell}(y),\quad k,\ell\in L_y.
\]
Equation \eqref{eq:mu-pentagon} shows that $\mu_y\in Z^2(L_y,\Ctimes)$.

\begin{theorem}\label{thm:simple-equivariantized-module}
As an abelian category,
\[
 \M(S,\Psi)^{L}
 \simeq
 \bigoplus_{\mathcal O\in L\backslash (\widehat{T}/S)}
 \Rep(L_y,\mu_y),
 \qquad y\in\mathcal O.
\]
In particular, $\M(S,\Psi)^L$ has rank one if and only if
\begin{enumerate}
\item $L$ acts transitively on $\widehat{T}/S$ through $U$; and
\item $\mu_y$ is nondegenerate on $L_y$ for one, equivalently every, $y\in \widehat{T}/S$.
\end{enumerate}
\end{theorem}

\begin{proof}
Apply Theorem~\ref{thm:orbit-stabilizer} to the action of $L$ on the simple
objects $\{\delta_y:y\in \widehat{T}/S\}$.  To spell this out, the action groupoid
$L\ltimes (\widehat{T}/S)$ has an arrow
\[
 [k,y]:y\longrightarrow U_k(y)
\]
for every $k\in L$ and $y\in \widehat{T}/S$.  Equation~\eqref{eq:mu-pentagon} is exactly the
groupoid cocycle identity \eqref{eq:groupoid-cocycle} for
\[
 \omega(k,\ell;y)=\mu_{k,\ell}(y).
\]
Thus the associated twisted algebra has basis $[k,y]$ and multiplication
\begin{equation}\label{eq:mu-twisted-groupoid-product}
 [k,U_\ell(y)]\,[\ell,y]
 =\mu_{k,\ell}(y)[k\ell,y],
\end{equation}
with product zero for noncomposable arrows.

For an $L$-orbit $\mathcal O\subseteq \widehat{T}/S$, an equivariant object supported on
$\mathcal O$ is a collection of vector spaces $V_y$, $y\in\mathcal O$, together
with maps
\[
 \rho(k,y):V_y\longrightarrow V_{U_k(y)}
\]
that satisfy the same relation as
\eqref{eq:mu-twisted-groupoid-product}.  It is therefore a module over the
corresponding orbit algebra.  If $y\in\mathcal O$ and $e_y=[1,y]$, then $e_y$ is
a full idempotent by Lemma~\ref{lem:transitive-groupoid-reduction}, and
\[
 e_y\,\mathbb C^{\mu}[L\ltimes\mathcal O]e_y
 =\bigoplus_{k\in L_y}\mathbb C[k,y]
 \cong\mathbb C^{\mu_y}[L_y].
\]
Consequently, the functor $N\mapsto e_yN$ gives an equivalence between the
objects supported on $\mathcal O$ and $\Rep(L_y,\mu_y)$.  Summing over the
orbits gives the displayed decomposition.  It has exactly one simple object if
and only if there is one orbit and the twisted stabilizer algebra is simple,
which is the stated rank-one criterion.
\end{proof}

\begin{example}[A transitive rank-one equivariantization]\label{ex:mixed-rank-one}
Let $G:=\Gm\times \mathbb{Z}/2\mathbb{Z}$.  Then $\widehat{T}=\mathbb{Z}$ and the action of $\mathbb{Z}/2\mathbb{Z}$ on $\widehat{T}$ is
trivial.  Take $S:=2\mathbb{Z}$, so $\widehat{T}/S\cong \mathbb{Z}/2\mathbb{Z}$, take $\Psi=1$, and let the nontrivial
element $g\in \mathbb{Z}/2\mathbb{Z}$ act by translation; that is,
\[
 U_g([m])=[m+1].
\]
With $\tau=1$ and $\mu=1$, all equations in Theorem~\ref{thm:explicit-equivariant-data} hold.  The action on $\widehat{T}/S$ is transitive and the
stabilizer is trivial, so $\M(2\mathbb{Z},1)^{\mathbb{Z}/2\mathbb{Z}}$ has rank one.  Notice that
$S\neq \widehat{T}$ and $U\neq\id$. \qed
\end{example}

\subsection{\texorpdfstring{Fiber functors and Hopf $2$-cocycles}{Fiber functors and Hopf 2-cocycles}}\label{sec:fiber-functors}

Call a sextuple
\[
 \mathcal D=(S,\Psi,L,U,\tau,\mu)
\]
{\bf admissible} if $S\subseteq \widehat{T}$ has finite index, $L\subseteq K$ is a subgroup, and the remaining data
satisfy Theorem~\ref{thm:explicit-equivariant-data}.  It is {\bf rank-one
admissible} if the two conditions in Theorem~\ref{thm:simple-equivariantized-module} hold.

Two admissible data are called {\bf equivalent} if, after conjugating one of the
subgroups $L$ by an element of $K$, the corresponding equivariant
$\Rep(T)$-module categories are equivalent.  This relation includes the usual
changes of cocycle representatives and of the scalar coherence cochains.

\begin{theorem}\label{thm:module-and-fiber-classification}
The following statements hold.
\begin{enumerate}
\item Equivalence classes of finite indecomposable semisimple $\Rep(G)$-module
categories are in bijection with equivalence classes of admissible data
\[
 (S,\Psi,L,U,\tau,\mu).
\]
\item Equivalence classes of fiber functors on $\Rep(G)$ are in bijection with
equivalence classes of rank-one admissible data.
\item Every fiber functor on $\Rep(G)$ is classical.
\end{enumerate}
\end{theorem}

\begin{proof}
(1) By Theorem~\ref{thm:tensor-equivalence}, Proposition~\ref{prop:equivariantization-modules}, and Proposition~\ref{prop:module-cats-torus}, every finite indecomposable module category has the
form
\[
 \bigl(\M(S,\Psi),U,\tau,\mu\bigr)^L,
\]
and the stated equivalence relation is precisely the one arising from those three
results.  

(2) This follows from Theorem~\ref{thm:simple-equivariantized-module}, because fiber functors are the same as
rank-one module categories.

(3) Let $\V$ be a rank-one module category corresponding to
rank-one admissible data, and let $V$ be its unique simple object.  Write the
underlying object of $V$ in $\M(S,\Psi)$ as
\[
 \overline V=\bigoplus_{y\in \widehat{T}/S}V_y\otimes\delta_y.
\]
Transitivity implies that every $V_y$ is nonzero and that all $\dim (V_y)$ are equal.
For an object $M=\bigoplus_yM_y\otimes\delta_y$ of $\M(S,\Psi)$, set
\[
 \ell(M):=\sum_{y\in \widehat{T}/S}\dim (M_y).
\]
If $X\in\Rep(G)$ and
$X|_T=\bigoplus_{x\in \widehat{T}}X_x\otimes\delta_x$, then the object-level action gives
\[
 \ell(X\otimes\overline V)
 =\left(\sum_{x\in \widehat{T}}\dim (X_x)\right)\ell(\overline V)
 =\dim(X)\ell(\overline V).
\]
Let $F_{\V}:\Rep(G)\to\Vect$ be the fiber functor defined by the action on the
rank-one category.  Since
\[
 X\otimes V\cong F_{\V}(X)\otimes V,
\]
the same length function gives
\[
 \dim (F_{\V}(X))\,\ell(\overline V)
 =\dim(X)\,\ell(\overline V).
\]
Hence $\dim (F_{\V}(X))=\dim (X)$ for every $X$.

Because $\Rep(G)$ is semisimple, an exact $\mathbb C$-linear functor to $\Vect$ is
determined up to natural isomorphism by its values on simple objects.  Choosing
vector-space isomorphisms on the simples therefore yields a natural isomorphism
between the underlying linear functor of $F_{\V}$ and the usual forgetful functor.
Thus $F_{\V}$ is classical.
\end{proof}

\begin{corollary}\label{cor:hopf-cocycle-classification}
Gauge-equivalence classes of Hopf $2$-cocycles on $\mathscr{O}(G)$ are in bijection
with equivalence classes of rank-one admissible data
\[
 (S,\Psi,L,U,\tau,\mu).
\]
\end{corollary}

\begin{proof}
Transport the tensor structure of a fiber functor along the natural isomorphism
with the forgetful functor supplied by Theorem~\ref{thm:module-and-fiber-classification}(3), and use
\S\ref{sec:hopf-cocycles}.  Monoidal natural isomorphisms correspond exactly to
gauge transformations.
\end{proof}

We now study the support and minimality of a Hopf $2$-cocycle on $\mathscr O(G)$.

\begin{theorem}\label{prop:factorization}
Let
\[
 \mathcal D:=(S,\Psi,L,U,\tau,\mu)
\]
be rank-one admissible data, and let $y_0:=S\in \widehat{T}/S$.  Since $S$ fixes
$y_0$, the formula
\begin{equation}\label{eq:toral-isotropy-cocycle}
 \psi_{\mathcal D}(s,t):=\Psi(s,t;y_0),
 \qquad s,t\in S,
\end{equation}
defines a normalized cocycle $\psi_{\mathcal D}\in Z^2(S,\Ctimes)$.  Let
\[
 R_{\mathcal D}(s,t)
 :=\frac{\psi_{\mathcal D}(s,t)}{\psi_{\mathcal D}(t,s)},
 \] 
 \[
 \Lambda_{\mathcal D}
 :=\Rad(R_{\mathcal D})
 =\{s\in S\mid R_{\mathcal D}(s,t)=1\text{ for all }t\in S\},
\]
and let
\[
 T_{\mathcal D}:=\Lambda_{\mathcal D}^{\perp}
 =\{a\in T\mid s(a)=1\text{ for every }s\in\Lambda_{\mathcal D}\}.
\]
Let $J_{\mathcal D}$ be the Hopf $2$-cocycle corresponding to $\mathcal D$, and let
$H_{\mathcal D}\subseteq G$ be its support.  Then after conjugating
$H_{\mathcal D}$ in $G$, one has
\begin{equation}\label{eq:support-from-admissible-data}
 \pi(H_{\mathcal D})=L,
 \qquad
 H_{\mathcal D}\cap T=T_{\mathcal D}.
\end{equation}
Consequently,
\begin{equation}\label{eq:minimality-full-data}
 J_{\mathcal D}\text{ is minimal}
 \quad\Longleftrightarrow\quad
 L=K\ \text{ and }\ \Lambda_{\mathcal D}=\{1\}.
\end{equation}
Equivalently, gauge-equivalence classes of minimal Hopf $2$-cocycles on
$\mathscr O(G)$ are parametrized by equivalence classes of admissible data
\[
 (S,\Psi,K,U,\tau,\mu)
\]
for which $K$ acts transitively on $\widehat{T}/S$, the stabilizer cocycle $\mu_{y_0}$ is
nondegenerate, and the alternating bicharacter $R_{\mathcal D}$ on $S$ is
nondegenerate.
\end{theorem}

\begin{proof}
Set $G_L:=\pi^{-1}(L)$.  The module category attached to $\mathcal D$ is first a
rank-one $\Rep(G_L)\simeq\Rep(T)^L$-module category, and its $\Rep(G)$-action is
obtained by restriction along $\Rep(G)\to\Rep(G_L)$.  Thus it is enough to compute
the support of the corresponding cocycle on $G_L$.

Let $V$ be the unique simple object of $\M(S,\Psi)^L$, and write its image in
$\M(S,\Psi)$ as
\[
 \overline V=\bigoplus_{y\in \widehat{T}/S}V_y\otimes\delta_y.
\]
By transitivity, every $V_y$ is nonzero and all the spaces $V_y$ have the same
dimension, say $d$.  Let $\mathcal A_{\mathcal D}$ be the internal endomorphism
algebra of $V$ in $\operatorname{Ind}(\Rep(G_L))$.  The induction--forgetful
adjunction for equivariantization identifies its image in
$\operatorname{Ind}(\Rep(T))$ with the internal endomorphism algebra
\[
 \mathcal B_{\mathcal D}
 :=\underline{\End}_{\M(S,\Psi)}(\overline V).
\]
Explicitly, $\mathcal B_{\mathcal D}$ is the $\widehat{T}$-graded algebra with
\[
 (\mathcal B_{\mathcal D})_x
 =\bigoplus_{y\in \widehat{T}/S}\Hom(V_y,V_{xy}),
 \qquad x\in \widehat{T},
\]
and, for $f\in(\mathcal B_{\mathcal D})_x$ and
$g\in(\mathcal B_{\mathcal D})_z$, multiplication is given by
\begin{equation}\label{eq:internal-end-product}
 (f\cdot g)_y
 =\Psi(x,z;y)\,f_{zy}\circ g_y.
\end{equation}

We now explain the groupoid reduction used to compute this algebra.  Let
\[
 \mathcal G:=A\ltimes (\widehat{T}/S).
\]
Its objects are the cosets $y\in \widehat{T}/S$, and for every $x\in \widehat{T}$ and $y\in \widehat{T}/S$ it
has an arrow
\[
 [x,y]:y\longrightarrow xy.
\]
The groupoid is transitive: if $y=aS$, then $[a,y_0]$ is an arrow from
$y_0=S$ to $y$.  Its isotropy group at $y_0$ is
\[
 \operatorname{Aut}_{\mathcal G}(y_0)
 =\{[x,y_0]\mid xy_0=y_0\}
 =\{[s,y_0]\mid s\in S\}\cong S.
\]
Moreover, the module-cocycle identity \eqref{eq:module-cocycle} is precisely the
groupoid cocycle identity \eqref{eq:groupoid-cocycle}.  We may therefore form the
twisted groupoid algebra
\[
 \mathcal G_{\Psi}:=\mathbb C^{\Psi}[\widehat{T}\ltimes (\widehat{T}/S)],
\]
whose multiplication is
\begin{equation}\label{eq:Psi-twisted-groupoid-product}
 [x,zy][z,y]=\Psi(x,z;y)[xz,y].
\end{equation}
After choosing bases $V_y\cong\mathbb C^d$, formula
\eqref{eq:internal-end-product} identifies $\mathcal B_{\mathcal D}$ with the
matrix-enriched version of this algebra:
\begin{equation}\label{eq:internal-end-as-groupoid-algebra}
 \mathcal B_{\mathcal D}
 \cong
 \operatorname{Mat}_d(\mathcal G_{\Psi}).
\end{equation}
Indeed, the copy of $\operatorname{Mat}_d(\mathbb C)$ attached to the arrow
$[x,y]$ is exactly $\Hom(V_y,V_{xy})$, and
\eqref{eq:Psi-twisted-groupoid-product} reproduces
\eqref{eq:internal-end-product}.

Let $e_y=[1,y]$.  The corner at the base point consists of the loops labelled by
$S$:
\begin{equation}\label{eq:isotropy-corner-S}
 e_{y_0}\mathcal G_{\Psi}e_{y_0}
 =\bigoplus_{s\in S}\mathbb C[s,y_0]
 \cong\mathbb C^{\psi_{\mathcal D}}[S],
\end{equation}
where the last assertion follows from
$[s,y_0][t,y_0]=\Psi(s,t;y_0)[st,y_0]$.
Choose representatives $r_y\in \widehat{T}$ with $r_yy_0=y$ and $r_{y_0}=1$, and let
$p_y,q_y$ be the normalized connecting arrows of
\eqref{eq:connecting-arrows}.  The explicit map
\[
 a\in e_{y'}\mathcal G_{\Psi}e_y
 \longmapsto
 E_{y',y}\otimes q_{y'}ap_y
\]
from \eqref{eq:explicit-groupoid-matrix-map} gives
\[
 \mathcal G_{\Psi}
 \cong
 \operatorname{Mat}_{|\widehat{T}/S|}\!\left(
   \mathbb C^{\psi_{\mathcal D}}[S]
 \right).
\]
Combining this with \eqref{eq:internal-end-as-groupoid-algebra} yields the desired
algebra isomorphism
\begin{equation}\label{eq:internal-end-matrix-algebra}
 \mathcal B_{\mathcal D}
 \cong
 \operatorname{Mat}_{d|\widehat{T}/S|}\!\left(\mathbb C^{\psi_{\mathcal D}}[S]\right).
\end{equation}

It remains to keep track of the $\widehat{T}$-grading and the center.  If $u_s$ denotes the
basis element of $\mathbb C^{\psi_{\mathcal D}}[S]$ corresponding to $s\in S$,
then the above matrix realization has the shifted grading
\begin{equation}\label{eq:shifted-matrix-grading}
 \deg(E_{y',y}\otimes u_s)=r_{y'}sr_y^{-1}.
\end{equation}
In particular, every diagonal element $E_{y,y}\otimes u_s$ has degree $s$,
because $\widehat{T}$ is abelian.  Furthermore,
\[
 u_su_t=\psi_{\mathcal D}(s,t)u_{st},
 \qquad
 u_tu_s=\psi_{\mathcal D}(t,s)u_{st},
\]
so
\begin{equation}\label{eq:center-of-twisted-S-algebra}
 Z\!\left(\mathbb C^{\psi_{\mathcal D}}[S]\right)
 =\bigoplus_{s\in\Lambda_{\mathcal D}}\mathbb C u_s.
\end{equation}
Since the center of a full matrix algebra consists of scalar matrices over the
center of its coefficient algebra,
\[
 Z(\mathcal B_{\mathcal D})
 =I_{d|\widehat{T}/S|}\otimes
  Z\!\left(\mathbb C^{\psi_{\mathcal D}}[S]\right).
\]
Under the inverse of the groupoid matrix map, the element
$I_{|\widehat{T}/S|}\otimes u_s$ is
\[
 \sum_{y\in \widehat{T}/S}p_yu_sq_y,
\]
a sum of loops, one at each object $y$.  Each summand has degree
$r_ysr_y^{-1}=s$.  Thus this central element is homogeneous of degree $s$, and
\eqref{eq:center-of-twisted-S-algebra} shows that the $T$-weights occurring in the
center are precisely the elements of $\Lambda_{\mathcal D}$.  In particular,
\begin{equation}\label{eq:center-weights}
 Z(\mathcal B_{\mathcal D})_x\neq0
 \quad\Longleftrightarrow\quad
 x\in\Lambda_{\mathcal D},
 \qquad
 \dim Z(\mathcal B_{\mathcal D})_1=1.
\end{equation}

By Theorem~\ref{thm:module-and-fiber-classification}, the fiber functor defined by
$\mathcal D$ is classical.  Therefore $\mathcal A_{\mathcal D}$ is the one-sided
twisted torsor algebra associated to the corresponding Hopf $2$-cocycle on
$G_L$ (up to passing to the opposite algebra, which does not change the center).
If $H_{\mathcal D}$ is its support, then
\cite[Theorem~3.5]{G1} gives an isomorphism of $G_L$-algebras
\begin{equation}\label{eq:center-support-homogeneous-space}
 Z(\mathcal A_{\mathcal D})\cong\mathscr O(G_L/H_{\mathcal D}).
\end{equation}
The degree-one part in \eqref{eq:center-weights} is one-dimensional.  On the other
hand, the $T$-orbits in $G_L/H_{\mathcal D}$ are indexed by
$L/\pi(H_{\mathcal D})$; hence
\[
 \dim\mathscr O(G_L/H_{\mathcal D})^T
 =[L:\pi(H_{\mathcal D})].
\]
Thus $\pi(H_{\mathcal D})=L$.

It follows that $T$ acts transitively on $G_L/H_{\mathcal D}$ and that
\[
 G_L/H_{\mathcal D}
 \cong T/(H_{\mathcal D}\cap T)
\]
as $T$-varieties.  Consequently, the group of $T$-weights occurring in
$\mathscr O(G_L/H_{\mathcal D})$ is
$(H_{\mathcal D}\cap T)^\perp$.  Comparing this with
\eqref{eq:center-weights} yields
\[
 (H_{\mathcal D}\cap T)^\perp=\Lambda_{\mathcal D},
\]
and therefore $H_{\mathcal D}\cap T=T_{\mathcal D}$.  This proves
\eqref{eq:support-from-admissible-data}.

Finally, the lift of the cocycle from $G_L$ to $G$ has the same support
$H_{\mathcal D}$.  Hence it is minimal on $G$ if and only if
$H_{\mathcal D}=G$, which, by \eqref{eq:support-from-admissible-data}, is equivalent
to $L=K$ and $\Lambda_{\mathcal D}=\{1\}$.  The last formulation simply combines
this condition with the rank-one criterion in
Theorem~\ref{thm:simple-equivariantized-module}.
\end{proof}

\begin{remark}\label{rem:two-nondegeneracy-conditions}
The two nondegeneracy conditions in Theorem~\ref{prop:factorization} play
different roles.  Nondegeneracy of $\mu_{y_0}$ on the finite stabilizer
$L_{y_0}$ is the rank-one condition, whereas nondegeneracy of
$R_{\mathcal D}$ on the lattice $S$ is the additional condition that the torus
part of the support be all of $T$.  In particular,
$\psi_{\mathcal D}$ is an isotropy cocycle extracted from the full module-category
data; it need not coincide with the restriction of $J_{\mathcal D}$ to
$\mathscr O(T)$, and it can encode mixed cocycle information. \qed
\end{remark}

\section{Commutative direct products and mixed cocycles}
\label{sec:commutative-direct-products}

Assume in this section that $G:=T\times K$ and that $K$ is finite {\bf abelian}.  Then $G$
is diagonalizable and
\[
 \mathscr{O}(G)=\mathbb{C}[\widehat T\times \widehat K].
\]
Hopf $2$-cocycles on $\mathscr{O}(G)$ are therefore ordinary normalized group $2$-cocycles on the
abelian group $\widehat T\times \widehat K$ \cite{G1}.

\subsection{The canonical K\"unneth decomposition}

\begin{theorem}\label{thm:kunneth}
There is a canonical group isomorphism
\begin{equation}\label{eq:kunneth}
 H^2(\widehat T\times \widehat K,\Ctimes)
 \cong
 H^2(\widehat T,\Ctimes)
 \times \Hom(\widehat T\otimes_{\mathbb Z}\widehat K,\Ctimes)
 \times H^2(\widehat K,\Ctimes).
\end{equation}
If $\psi$ and $\xi$ are normalized $2$-cocycles on $\widehat T$ and $\widehat K$, respectively, and
$\chi:\widehat T\times \widehat K\to\Ctimes$ is a bicharacter, then the corresponding class is
represented by
\begin{equation}\label{eq:mixed-cocycle}
 J_{\psi,\chi,\xi}\bigl((x,\alpha),(y,\beta)\bigr)
 =\psi(x,y)\chi(x,\beta)\xi(\alpha,\beta).
\end{equation}
The inverse of \eqref{eq:kunneth} sends $[J]$ to
\[
 \left(
 [J|_{(\widehat T\times\{1\})^2}],
 \chi_J,
 [J|_{(\{1\}\times \widehat K)^2}]
 \right),
 \qquad
 \chi_J(x,\beta)=R^J((x,1),(1,\beta)).
\]
\end{theorem}

\begin{proof}
Because $\Ctimes$ is divisible, it is an injective $\mathbb Z$-module.  The
universal coefficient sequence therefore identifies
\[
 H^2(A,\Ctimes)\cong\Hom(\textstyle\bigwedge^2 A,\Ctimes)
\]
for every abelian group $A$; under this identification, a cocycle class is sent to
its alternating bicharacter $R^J:=J_{21}^{-1}J$.  The canonical decomposition
\[
 \bigwedge^2(\widehat T\times \widehat K)
 \cong
 \bigwedge^2\widehat T\oplus(\widehat T\otimes_{\mathbb Z}\widehat K)\oplus\bigwedge^2\widehat K
\]
then gives \eqref{eq:kunneth}.  Formula \eqref{eq:mixed-cocycle} is a normalized
$2$-cocycle, and its alternating bicharacter has the prescribed three components.
This also proves the formula for the inverse.
\end{proof}

Equivalently, a bicharacter $\chi$ is the same as a homomorphism
\[
 \lambda:\widehat T\longrightarrow K=\widehat {\widehat K},
 \qquad
 \chi(x,\beta)=\beta(\lambda(x)).
\]
In these terms,
\begin{equation}\label{eq:R-block}
 R^J((x,\alpha),(y,\beta))
 =R^\psi(x,y)\,
 \beta(\lambda(x))\,
 \alpha(\lambda(y))^{-1}\,
 R^\xi(\alpha,\beta).
\end{equation}

\subsection{Minimality}

\begin{proposition}\label{prop:radical}
For the cocycle in \eqref{eq:mixed-cocycle}, written by means of $\lambda$, one has
\[
 \Rad(R^J)=
 \left\{
 (x,\alpha)\in \widehat T\times \widehat K\ \middle|\
 \begin{array}{l}
 R^\psi(x,-)=\alpha\circ\lambda,\\[1mm]
 R^\xi(\alpha,-)=\operatorname{ev}_{\lambda(x)}^{-1}
 \end{array}
 \right\}.
\]
Consequently, $J$ is minimal if and only if this set is trivial.
\end{proposition}

\begin{proof}
Fix $(x,\alpha)\in \widehat T\times \widehat K$. Then setting $\beta=1$ in
\eqref{eq:R-block} shows that 
$$R^J((x,\alpha),(y,1))=1$$
for every $y$ exactly when
$R^\psi(x,-)=\alpha\circ\lambda$.  Setting $y=1$ in
\eqref{eq:R-block} shows that
$$R^J((x,\alpha),(1,\beta))=1$$
for every $\beta$ exactly when
$R^\xi(\alpha,-)=\operatorname{ev}_{\lambda(x)}^{-1}$.  Since the two factors
generate $\widehat T\times \widehat K$, these conditions are necessary and sufficient.
\end{proof}

\subsection{Comparison with the equivariantization classification}
\label{subsec:kunneth-versus-equivariantization}

The classification in Theorem~\ref{thm:kunneth} and the classification in
Corollary~\ref{cor:hopf-cocycle-classification} describe the same
gauge-equivalence classes, but they use rather different coordinates.  The
K\"unneth coordinates are canonical and cohomological, whereas the data of
\S\ref{sec:fiber-functors} are orbit--stabilizer data arising from
equivariantization.  We now make the relation between them explicit.

Let
\[
 [J]=[J_{\psi,\chi,\xi}]
 \in H^2(\widehat T\times \widehat K,\Ctimes)
\]
be written in the K\"unneth form of Theorem~\ref{thm:kunneth}, and let
\[
 \lambda:\widehat T\longrightarrow K=\widehat{\widehat K},
 \qquad
 \chi(x,\beta)=\beta(\lambda(x)),
\]
be the homomorphism corresponding to the mixed bicharacter.  Define
homomorphisms
\[
 r_\psi:\widehat T\longrightarrow T,
 \qquad
 y(r_\psi(x))=R^\psi(x,y),
\]
and
\[
 r_\xi:\widehat K\longrightarrow K,
 \qquad
 \beta(r_\xi(\alpha))=R^\xi(\alpha,\beta).
\]
Also let
\[
 \lambda^\vee:\widehat K\longrightarrow T,
 \qquad
 x(\lambda^\vee(\alpha))=\alpha(\lambda(x)),
\]
be the homomorphism dual to $\lambda$.  Let
\[
 T_\psi:=\overline{\operatorname{im}(r_\psi)}
       =\Rad(R^\psi)^\perp\subseteq T,
 \qquad
 N_\xi:=\operatorname{im}(r_\xi)
       =\Rad(R^\xi)^\perp\subseteq K.
\]
Thus, $T_\psi$ and $N_\xi$ are the supports of the two diagonal K\"unneth
components.

We have the following dictionary between the two classifications.

\begin{proposition}
\label{prop:kunneth-equivariantization-dictionary}
Let $\mathcal D_J$ be the equivalence class of rank-one admissible data in
\S\ref{sec:fiber-functors} corresponding to $J$.  Then one may choose a
representative
\[
 \mathcal D_J=(S_J,\Psi_J,L_J,U_J,\tau_J,\mu_J)
\]
with the following properties:
\begin{enumerate}

\item The subgroup and orbit data are
\begin{equation}\label{eq:dictionary-S-L}
 N:=N_\xi,
 \qquad
 L_J=\lambda(\widehat T)N,
 \qquad
 S_J=\lambda^{-1}(N).
\end{equation}
There is a canonical isomorphism
\begin{equation}\label{eq:dictionary-orbit}
 \widehat T/S_J\xrightarrow{\cong}L_J/N,
 \qquad
 xS_J\longmapsto\lambda(x)N,
\end{equation}
and, under this identification, $U_J$ is the left translation action of
$L_J$ on $L_J/N$.  In particular, the stabilizer of the base point is $N$.

\item The restriction class $[\xi]\in H^2(\widehat K,\Ctimes)$ determines the
nondegenerate stabilizer class
\[
 [\mu_J]\in H^2(N,\Ctimes).
\]
More precisely, if $B=\Rad(R^\xi)$, then
$\widehat K/B\cong\widehat N$, and the alternating form induced by $R^\xi$ on
$\widehat K/B$ is nondegenerate.  After choosing a simple block of
$\mathbb C^\xi[B]$, the stabilizer $N$ acts projectively on its unique simple
module; the multiplier of this projective action is $\mu_J$, and it is
nondegenerate.  The resulting cohomology class $[\mu_J]$ depends only on
$[\xi]$, up to the harmless inversion dictated by the convention for
projective representations.

The cochains $\Psi_J$ and $\tau_J$ encode the toral and mixed components
relative to this finite block; hence they can also involve $\xi$ through the
chosen projective intertwiners.  They depend on choices of cocycle
representatives and intertwiners, but their equivalence class does not.  In
particular, the isotropy cocycle
\[
 \psi_{\mathcal D_J}(s,t)=\Psi_J(s,t;S_J)
\]
need not coincide with the restriction
$\psi|_{S_J\times S_J}$; it is the Mackey obstruction obtained after the
mixed and finite data have been incorporated.

\item The K\"unneth coordinates are recovered from the cocycle associated
with $\mathcal D_J$ by
\begin{equation}\label{eq:recover-kunneth}
 \left[
 J\big|_{(\widehat T\times\{1\})^2}
 \right]=[\psi],
 \quad
 R^J((x,1),(1,\beta))=\chi(x,\beta),
 \quad
 \left[
 J\big|_{(\{1\}\times \widehat K)^2}
 \right]=[\xi].
\end{equation}
Consequently, two rank-one admissible data give the same K\"unneth triple if
and only if they are equivalent in the sense of
\S\ref{sec:fiber-functors}.

\item Let $H_J\leq T\times K$ be the support of $J$, and let
$\Lambda_{\mathcal D_J}$ be the isotropy radical from
Proposition~\ref{prop:factorization}.  Then
\begin{equation}\label{eq:dictionary-support}
 \pi(H_J)=L_J=\lambda(\widehat T)N_\xi
\end{equation}
and
\begin{equation}\label{eq:dictionary-isotropy-radical}
 \Lambda_{\mathcal D_J}
 =\operatorname{pr}_{\widehat T}\Rad(R^J)
 =
 \left\{
 x\in \widehat T\ \middle|\
 \begin{array}{l}
 \text{there exists }\alpha\in \widehat K\text{ such that}\\[1mm]
 R^\psi(x,-)=\alpha\circ\lambda,\,
 R^\xi(\alpha,-)=\operatorname{ev}_{\lambda(x)}^{-1}
 \end{array}
 \right\}.
\end{equation}

\end{enumerate}
\end{proposition}

\begin{proof}
Consider first the twisted group algebra $\mathbb C^\xi[\widehat K]$, and let
\[
 B:=\Rad(R^\xi).
\]
Since the restriction of $\xi$ to $B$ has trivial alternating bicharacter,
it is a coboundary; after replacing $\xi$ by a cohomologous cocycle, we may
therefore assume that this restriction is trivial.  Then
\[
 Z(\mathbb C^\xi[\widehat K])=\mathbb C[B],
\]
and the simple matrix blocks of $\mathbb C^\xi[\widehat K]$ are indexed by
\[
 \widehat{B}\cong K/N_\xi.
\]

For $x\in \widehat T$, tensoring a projective $\widehat K$-module by the character
$\chi(x,-)$ translates its block by $\lambda(x)N_\xi$.  Hence the orbit of
the base block is
\[
 \lambda(\widehat T)N_\xi/N_\xi=L_J/N_\xi,
\]
and its stabilizer in $\widehat T$ is $\lambda^{-1}(N_\xi)$.  This gives
\eqref{eq:dictionary-S-L}, \eqref{eq:dictionary-orbit}, and the translation
action $U_J$.

The alternating form induced by $R^\xi$ on $\widehat K/B$ is nondegenerate.  The
natural action of the stabilizer $N_\xi$ on a chosen simple module $V$ in
the base block is projective; let $\mu_J$ be its multiplier.  Since
\[
 (\dim V)^2=|\widehat K/B|=|N_\xi|,
\]
the induced surjection
\[
 \mathbb C^{\mu_J}[N_\xi]\twoheadrightarrow\End(V)
\]
is an isomorphism.  Hence $\mu_J$ is nondegenerate.  Choosing compatible
intertwiners then produces $\Psi_J$, $\tau_J$, and $\mu_J$; changing these
choices changes only the representative of the admissible data.
Formula~\eqref{eq:recover-kunneth} is exactly the inverse map in
Theorem~\ref{thm:kunneth}.

For the support statement, define
\begin{equation}\label{eq:rho-J}
 \rho_J:\widehat T\times \widehat K\longrightarrow T\times K,
 \qquad
 \rho_J(x,\alpha)
 =
 \bigl(
 r_\psi(x)\lambda^\vee(\alpha)^{-1},
 \lambda(x)r_\xi(\alpha)
 \bigr).
\end{equation}
Using \eqref{eq:R-block}, one obtains
\[
 R^J((x,\alpha),(y,\beta))
 =(y,\beta)\bigl(\rho_J(x,\alpha)\bigr).
\]
Therefore
\[
 \Rad(R^J)=\ker(\rho_J),
 \qquad
 H_J=\overline{\operatorname{im}(\rho_J)}.
\]
Taking the second component in \eqref{eq:rho-J} gives
\[
 \pi(H_J)
 =
 \lambda(\widehat T)\operatorname{im}(r_\xi)
 =
 \lambda(\widehat T)N_\xi,
\]
which proves \eqref{eq:dictionary-support}.

Since
\[
 H_J^\perp=\Rad(R^J),
\]
elementary character duality gives
\[
 (H_J\cap T)^\perp=\operatorname{pr}_{\widehat T}\Rad(R^J).
\]
Indeed, a character of $T$ trivial on $H_J\cap T$ descends to a character of
\[
 H_J/(H_J\cap T)\subseteq K,
\]
and every character of a subgroup of the finite abelian group $K$ extends
to $K$.  Proposition~\ref{prop:factorization} now gives
\[
 \Lambda_{\mathcal D_J}
 =(H_J\cap T)^\perp
 =\operatorname{pr}_{\widehat T}\Rad(R^J).
\]
The explicit description in
\eqref{eq:dictionary-isotropy-radical} follows from
Proposition~\ref{prop:radical}.
\end{proof}

The preceding dictionary makes the comparison of the two minimality criteria
completely transparent.

\begin{corollary}
\label{cor:kunneth-equivariantization-minimality}
With the notation above, the following conditions are equivalent:
\begin{enumerate}

\item $J$ is minimal.

\item
$
 \Rad(R^J)=\{1\}$.

\item
$
 \Rad(R^\psi)=\{1\}$ and 
$
 \ker(\lambda^\vee)\cap\Rad(R^\xi)=\{1\}$.

\item
$
 T_\psi=T$ and $K=\lambda(\widehat T)N_\xi$.

\item the \S\ref{sec:fiber-functors} datum satisfies
$
 L_J=K$ and $\Lambda_{\mathcal D_J}=\{1\}$.
\end{enumerate}
Thus condition~(5) is precisely the minimality criterion of
Proposition~\ref{prop:factorization}, while conditions~(2)--(4) are its
K\"unneth formulations.
\end{corollary}

\begin{proof}
The equivalence of (1) and (2) is the definition of minimality for a
diagonalizable group.  

If (2) holds, then testing the radical on $(x,1)$ and
$(1,\alpha)$ gives
\[
 \ker(\lambda)\cap\Rad(R^\psi)=\{1\},
 \qquad
 \ker(\lambda^\vee)\cap\Rad(R^\xi)=\{1\}.
\]
Since $\lambda(\widehat T)$ is finite, $\ker(\lambda)$ has finite index in the
torsion-free group $\widehat T$.  Hence every nontrivial subgroup of $\widehat T$ meets
$\ker(\lambda)$ nontrivially, and therefore
\[
 \ker(\lambda)\cap\Rad(R^\psi)=\{1\}
 \quad\Longleftrightarrow\quad
 \Rad(R^\psi)=\{1\}.
\]
This proves $(2)\Rightarrow(3)$.

Conversely, assume (3), and let
\[
 (x,\alpha)\in\Rad(R^J).
\]
If $m$ is the order of $\alpha$, then
\[
 (x^m,1)\in\Rad(R^J)
\]
because $R^J$ is a bicharacter.  Hence
\[
 x^m\in\Rad(R^\psi)=\{1\}.
\]
Since $\widehat T$ is torsion-free, $x=1$.  It follows that
\[
 \alpha\in
 \ker(\lambda^\vee)\cap\Rad(R^\xi)
 =\{1\},
\]
so (2) holds.

By annihilator duality,
\[
 \Rad(R^\psi)^\perp=T_\psi,
 \qquad
 \ker(\lambda^\vee)^\perp=\lambda(\widehat T),
 \qquad
 \Rad(R^\xi)^\perp=N_\xi.
\]
Therefore (3) and (4) are equivalent.  

Finally,
\eqref{eq:dictionary-support} gives
\[
 L_J=K
 \quad\Longleftrightarrow\quad
 \ker(\lambda^\vee)\cap\Rad(R^\xi)=\{1\}.
\]
Moreover, by \eqref{eq:dictionary-isotropy-radical},
\[
 \Lambda_{\mathcal D_J}=\{1\}
 \quad\Longleftrightarrow\quad
 \Rad(R^\psi)=\{1\}.
\]
Indeed,
\[
 \Rad(R^\psi)\cap\ker(\lambda)
 \subseteq
 \operatorname{pr}_{\widehat T}\Rad(R^J).
\]
Conversely, if
\[
 x\in\operatorname{pr}_{\widehat T}\Rad(R^J),
\]
choose $\alpha$ with $(x,\alpha)\in\Rad(R^J)$ and raise to the order of
$\alpha$.  A power of $x$ then lies in
\[
 \Rad(R^\psi)\cap\ker(\lambda).
\]
Since $\widehat T$ is torsion-free and $\ker(\lambda)$ has finite index, the asserted
equivalence follows.  This proves the equivalence with (5).
\end{proof}

\begin{remark}
\label{rem:mixed-component-comparison}
The subgroup $N_\xi$ (and not $L_J$) is the subgroup carrying the
nondegenerate finite stabilizer cocycle $\mu_J$.  The mixed component enlarges
it to
\[
 L_J=\lambda(\widehat T)N_\xi,
\]
and the quotient
\[
 L_J/N_\xi\cong \widehat T/S_J
\]
is exactly the transitive orbit appearing in the rank-one criterion of
\S\ref{sec:fiber-functors}.  Thus, even for a minimal cocycle, one may
have
\[
 N_\xi\subsetneq K
\]
and a degenerate restriction $\xi$ on $\widehat K$; the missing finite directions are
supplied by $\lambda(\widehat T)$.

On the torus side the situation is different: since $\widehat T$ is torsion-free and
$\widehat K$ is finite, a nonzero element of $\Rad(R^\psi)$ always has a nonzero power
in $\ker(\lambda)$ and therefore gives a nontrivial element of
$\Rad(R^J)$.  Hence minimality forces $R^\psi$ itself to be nondegenerate.
The mixed component can compensate for degeneracy of the finite block, but
not for degeneracy of the toral block.

If the mixed component is trivial, then $\lambda=1$, so
\[
 S_J=\widehat T,
 \qquad
 L_J=N_\xi,
 \qquad
 U_J=\id.
\]
In this special case minimality does reduce to separate nondegeneracy of the
two diagonal blocks.  By contrast, the image of the mixed component in
$K/N_\xi$ is manifested in \S\ref{sec:fiber-functors} by the quotient
$\widehat T/S_J$ and its transitive $L_J$-action.  The part of $\lambda(\widehat T)$ lying
inside $N_\xi$ does not enlarge the orbit; instead, it is absorbed into the
cochains $\Psi_J$ and $\tau_J$ and can change the isotropy cocycle even when
$S_J=\widehat T$. \qed
\end{remark}

\begin{remark}
The two classifications have complementary advantages.  The K\"unneth
classification is canonical, carries the abelian group structure on
\[
 H^2(\widehat T\times \widehat K,\Ctimes)
\]
componentwise, and displays the mixed term explicitly.  The
equivariantization classification is less canonical at the cochain level,
but it records the orbit, stabilizer, support, and factorization of the
associated module category, and it continues to make sense for nonsplit
extensions and for noncommutative component groups.

Notice also that the cochain $\tau_J$ in the admissible datum is not, in
general, a canonical homomorphism $K\to T$.  The canonical mixed invariant is
\[
 \lambda:\widehat T\to K,
\]
or equivalently its dual
\[
 \lambda^\vee:\widehat K\to T.
\]
Converting this into a map $K\to T$ would require a noncanonical
identification $K\cong\widehat K$.  In the commutative direct-product case,
the identities above show that the two classifications agree exactly. \qed
\end{remark}

\begin{example}[A purely mixed class]\label{ex:purely-mixed}
Let $G:=\Gm\times \mathbb{Z}/2\mathbb{Z}$, so $\widehat T=\mathbb Z$ and $\widehat K=\mathbb{Z}/2\mathbb{Z}$.  Writing elements as
$(m,\varepsilon)$ with $\varepsilon\in\{0,1\}$, define
\[
 J((m,\varepsilon),(n,\delta))=(-1)^{m\delta}.
\]
Both restrictions of $J$ to the two factors are trivial, but
\[
 R^J((1,0),(0,1))=-1.
\]
Thus $[J]$ is the nontrivial mixed class in
$\Hom(\mathbb Z\otimes (\mathbb{Z}/2\mathbb{Z}),\Ctimes)$.  It is invisible on both diagonal factors and is detected entirely by the mixed component. \qed
\end{example}

\begin{example}[A minimal cocycle with degenerate finite block]
Let
$$
 G:=(\Gm)^2\times (\mathbb{Z}/2\mathbb{Z}),$$
 so that
 $$\widehat T=\mathbb Z^2,
 \qquad
 \widehat K=\mathbb{Z}/2\mathbb{Z},$$
and choose $q\in\Ctimes$ that is not a root of unity.  Define
\[
 J\bigl(((a,b),\varepsilon),((c,d),\delta)\bigr)
 =q^{ad}(-1)^{a\delta}.
\]
Then
\[
 R^J\bigl(((a,b),\varepsilon),((c,d),\delta)\bigr)
 =q^{ad-cb}(-1)^{a\delta+c\varepsilon}.
\]
If $((a,b),\varepsilon)$ lies in the radical, testing against $((0,1),0)$ gives
$q^a=1$, hence $a=0$.  Testing against $((1,0),0)$ then gives
$q^{-b}(-1)^\varepsilon=1$, hence $b=0$ and $\varepsilon=0$.  Therefore $R^J$ is
nondegenerate, although the restriction of $J$ to $\mathbb{Z}/2\mathbb{Z}\times \mathbb{Z}/2\mathbb{Z}$ is trivial.
This illustrates why minimality must be tested on the full block bicharacter. \qed
\end{example}

%\enlargethispage{8\baselineskip}
\begin{thebibliography}{ABCD}

\bibitem[A]{A}
E. Abe, \emph{Hopf algebras}, Cambridge Tracts in Mathematics, vol. 74,
Cambridge University Press, Cambridge--New York, 1980.

\bibitem[DM]{DM}
P. Deligne and J. Milne, Tannakian categories, in \emph{Hodge Cycles, Motives, and
Shimura Varieties}, Lecture Notes in Mathematics, vol. 900, Springer, 1982,
pp. 101--228.

\bibitem[EG1]{EG1}
P. Etingof and S. Gelaki, On cotriangular Hopf algebras,
\emph{Amer. J. Math.} \textbf{123} (2001), 699--713.

\bibitem[EG2]{EG2}
P. Etingof and S. Gelaki, Quasisymmetric and unipotent tensor categories,
\emph{Math. Res. Lett.} \textbf{15} (2008), no. 5, 857--866.

\bibitem[ENO]{ENO}
P. Etingof, D. Nikshych, and V. Ostrik, Weakly group-theoretical and solvable
fusion categories, \emph{Adv. Math.} \textbf{226} (2011), 176--205.

\bibitem[EGNO]{EGNO}
P. Etingof, S. Gelaki, D. Nikshych, and V. Ostrik,
\emph{Tensor Categories}, Mathematical Surveys and Monographs, vol. 205,
American Mathematical Society, Providence, RI, 2015.

\bibitem[G1]{G1}
S. Gelaki, Twisting of affine algebraic groups, I,
\emph{Int. Math. Res. Not. IMRN} (2015), no. 16, 7552--7574.

\bibitem[G2]{G2}
S. Gelaki, Module categories over affine group schemes,
\emph{Quantum Topol.} \textbf{6} (2015), no. 1, 1--37.

\bibitem[Mon]{M}
S. Montgomery, \emph{Hopf Algebras and Their Actions on Rings},
CBMS Regional Conference Series in Mathematics, vol. 82,
American Mathematical Society, Providence, RI, 1993.

\bibitem[Mov]{Mo}
M. V. Movshev, Twisting in group algebras of finite groups,
\emph{Funct. Anal. Appl.} \textbf{27} (1993), no. 4, 240--244.

\bibitem[S]{S}
T. A. Springer, \emph{Linear Algebraic Groups}, second ed.,
Progress in Mathematics, vol. 9, Birkh\"auser, Boston, 1998.

\end{thebibliography}

\end{document}
